\section{Experimental Design}
\label{sec:design}

The measurement protocol and the decision rule are fixed before the headline
comparison, so the evaluation year cannot shape the design judged on it. That
ordering matters most for the frozen comparator, since a static width chosen
after seeing 2024 would flatter any policy measured against it. The interface
corrections that preceded the final runs are reported as corrections rather
than absorbed into the design, because they changed what the policy could
observe and what it was rewarded for. Table~\ref{tab:protocol} states the
pools, the split and the training budget, and Fig.~\ref{fig:split} shows how
the windows sit on the calendar.

\begin{table}[tbp]
\caption{The study protocol and the training configuration. Both pools are
USDC/WETH with USDC as token0; prices are quoted in USDC per WETH.}
\label{tab:protocol}
\centering
\small
\begin{tabular}{@{}L{0.34\linewidth}L{0.60\linewidth}@{}}
\toprule
\textbf{Setting} & \textbf{Value} \\
\midrule
Primary pool & 0.05\,\% (500 pips), tick spacing 10 \\
Contrast pool & 0.30\,\% (3000 pips), tick spacing 60 \\
Training window & 2022-01-01 to 2023-12-31 UTC \\
Evaluation window & 2024-01-01 to 2024-12-31 UTC \\
Episode & thirty days, one decision every ten minutes \\
Capital & \usdc{100000} per episode \\
Evaluation tiling & 12 windows of 30 days plus a partial final window (13 episodes) \\
Seeds & 5 \\
Training budget & 1\,000\,000 timesteps per seed \\
Timing & one deployment per episode; the band is held afterwards \\
Width arms & half-widths of 50, 250 and 1000 tick spacings, plus hold \\
Hedge head & short ratio in $\{0,0.25,0.5,0.75,1.0\}$ at every decision \\
Risk penalty & CVaR at $5\,\%$ of the trailing one-day loss, $\lambda=0.05$ \\
Checkpoint selection & highest validation Sharpe on the last 180 days of training \\
Algorithm & clipped PPO, $\gamma=0.9999$, GAE $\lambda=0.95$, clip $0.15$ \\
Network & 2 hidden layers of 256 units \\
Parallel environments & 4 \\
\bottomrule
\end{tabular}
\end{table}

\begin{figure}[tbp]
\centering
\begin{tikzpicture}[font=\small, x=1.4mm, y=2.4mm]
  \draw[line width=3.2pt, black!25] (0,1.6) -- (66,1.6);
  \draw[line width=3.2pt, black!55] (66,1.6) -- (100,1.6);
  \foreach \i in {0,...,12} {
    \draw[black!70] (66+\i*2.6,1.1) -- (66+\i*2.6,2.1);
  }
  \draw[-{Stealth[length=2.2mm]}, thick] (0,0.5) -- (104,0.5);
  \foreach \x/\lab in {0/2022, 33/2023, 66/2024, 99/2025} {
    \node[below, font=\scriptsize] at (\x,0.5) {\lab};
  }
  \node[above, font=\scriptsize] at (33,2.0) {training: windows of 30 days resampled};
  \node[above, font=\scriptsize] at (84,2.0) {held out};
  \draw[dashed, black!60] (66,0.5) -- (66,3.0);
  \node[above, font=\scriptsize, black!60] at (66,3.0) {split};
\end{tikzpicture}
\caption{The walk-forward split. Training episodes are windows of thirty days
sampled from 2022--2023; evaluation tiles 2024 into windows of thirty days, the
last a partial period, and no episode crosses the split.}
\label{fig:split}
\end{figure}

\subsection{Benchmarks}

A learned policy that beats only the weakest rule written by hand has
demonstrated nothing, so the ladder in Table~\ref{tab:benchmarks} includes the
opponents that could otherwise account for a positive result. It runs from the
one-shot rules through the rebalancing heuristics to the best static width,
chosen on the training split and frozen there, and it ends with a delta-matched
hedge. That last rung matters because a policy whose advantage is the removal
of directional risk would otherwise be measured against opponents that carry
that risk. Two reference holdings frame the ladder. A pure buy-and-hold basket
holds the opening token mix as tokens and never provides liquidity, which makes
it the numeraire for excess over HODL. The \textsf{hodl} rule deploys the
opening band before the first decision and never recentres, so it earns fees
and pays impermanent loss.

\begin{table}[tbp]
\caption{Benchmark policies. The learned policy is compared with all of them;
a positive claim requires beating the best static width and the delta-matched
hedged benchmark, not just \textsf{hodl}.}
\label{tab:benchmarks}
\centering
\small
\begin{tabular}{@{}L{0.28\linewidth}L{0.64\linewidth}@{}}
\toprule
\textbf{Benchmark} & \textbf{Rule} \\
\midrule
\textsf{hodl} & Hold the opening band of half-width 120 ticks; never recentre. \\
\textsf{passive\_narrow} & Deploy a $\pm 5\,\%$ band once; never recentre. \\
\textsf{passive\_wide} & Deploy a $\pm 20\,\%$ band once; never recentre. \\
\textsf{full\_range\_v2} & The Uniswap-v2-equivalent full range. \\
\textsf{tau\_reset} & Recentre whenever the price exits the live band. \\
\textsf{cost\_aware} & Recentre at most daily, and only when the estimated fee gain exceeds gas. \\
Best static width & The ladder width with the best Sharpe on the training split, frozen for test. \\
Hedged HODL & The opening band held, plus a short perpetual matched to its delta, with funding and perp costs. \\
Learned policy & PPO over the joint (timing, width, hedge ratio) action space. \\
\bottomrule
\end{tabular}
\end{table}

\subsection{Metrics}

Two primary statistics carry the inference, both computed on the joint equity
curve of the liquidity and hedge legs, so a hedge that flattens the tail at a
cost in return is scored on the trade rather than on one side of it. The
annualised Sharpe ratio is the mean of the per-step equity returns divided by
their standard deviation, scaled to a year at the ten-minute cadence, and it
measures the return per unit of risk. CVaR at the $5\,\%$ level is the average
of the worst $5\,\%$ of per-step equity changes in USDC, the same quantity the
reward penalty controls, and it measures the depth of the loss when the market
turns. The learned policy is compared with the benchmark ladder through paired
per-window differences in both statistics, with bootstrap confidence intervals
over the thirteen windows. Because the tail is thin at that sample size, the
worst window is reported alongside the intervals. Secondary statistics are raw
PnL, excess over HODL, annualised return, Sortino ratio, maximum drawdown,
rebalance count, and the cost ledger of fee income, gas, slippage, realised
impermanent loss and the hedge flows. Per-regime tables stratify the windows
by their dominant label, namely bull, bear, sideways or high volatility.

\subsection{Ablations and the boundary}

The failure modes were repaired in sequence rather than varied in a factorial
design, and the sequence is reported as such. Table~\ref{tab:ablations} lists
the design choices the repairs touched, namely band geometry, the reward basket
and window sampling, the rebalance interval, the hedge head and the risk
penalty. A sequence cannot isolate the marginal effect of a single change,
because a repair landed on an interface the next repair could also move, so
Section~\ref{sec:results} reports the progression rather than an attribution.
Three targeted probes complete the table and recover part of what the sequence
cannot separate. The friction terms are ablated on a one-shot band and on a
recentring rule, the rule ladder is repeated on the $0.30\,\%$ tier, and an arm
with the regime feature frozen is trained across the five seeds.

The boundary question calls for a sweep of realised volatility, volume
relative to liquidity, fee tier, gas price and capital, evaluated on replayed
windows rather than a fitted process. The choice keeps a boundary point tied to
market states that occurred, so the sweep does not add a generator's
assumptions to those already carried by the testbed. The completed sweep moves
the gas and capital levels and the fee tier, while the realised windows supply
the volatility and volume variation the LVR overlay uses. A controlled
volatility and volume sweep remains open.

\begin{table}[tbp]
\caption{The design choices the repair sequence touched, plus the targeted
probes that complete the table. The repair changes were made in sequence rather
than varied one at a time, so Section~\ref{sec:results} measures the sequence
of repairs, not the marginal effect of each change; the friction, fee-tier and
regime-feature rows are the probes.}
\label{tab:ablations}
\centering
\small
\begin{tabular}{@{}L{0.24\linewidth}L{0.34\linewidth}L{0.34\linewidth}@{}}
\toprule
\textbf{Design choice} & \textbf{Ablation} & \textbf{Question it answers} \\
\midrule
Band geometry & removed from the observation & does the agent need to see the live band to time rebalances? \\
Window sampling & windows held fixed instead of resampled & how much does resampling the training windows buy? \\
Hedge & overlay disabled & what does the hedge contribute net of funding? \\
Risk penalty & toggled & does the risk penalty change the tail without costing Sharpe? \\
Regime features & frozen to a constant across the five seeds & do the regime labels inform the policy? \\
Frictions & full model, no gas, no slippage, flat gas, no impermanent loss & which cost drives the result? \\
Fee tier & repeated on the $0.30\,\%$ pool & does a higher fee tier buy a longer viable range? \\
\bottomrule
\end{tabular}
\end{table}

The pre-registered rule awards a positive verdict only if the learned policy's
mean Sharpe beats the best static width and the delta-matched hedge with a
paired $95\,\%$ confidence interval that excludes zero, and only if the
policy's CVaR is no worse than the best heuristic's. Both clauses must hold,
so tail protection without a return advantage does not pass.
Section~\ref{sec:results} reports where the rule is met and where it is not.
Where the evidence is negative, the negative result is the finding.
